\chapter{Evaluation}
\label{ch:evaluation}

The Introduction advanced two claims: that zero-configuration protocols can provision AI services (\argref{C1}), and that network-level provisioning reduces total configuration effort (\argref{C2}). This chapter evaluates both. I assess \argref{C1} through an implementation census and protocol analysis, and \argref{C2} through a cognitive walkthrough following the methodology of Wharton et al.~\cite{wharton1994walkthrough}. The security trade-offs of broadcast discovery and a usability critique against Nielsen's heuristics~\cite{nielsen1994heuristics} receive their treatment in Chapter~\ref{ch:discussion}. The chapter closes with threats to validity. Both evaluations are analytical rather than empirical --- Saturn has no user base, no production deployment, and no institutional partner willing to run a field study --- and analytical methods are appropriate for a system at this stage because they evaluate design properties without requiring participants.

%----------------------------------------------------------------------
\section{C1: Zero-Configuration AI Provisioning Is Feasible}
\label{sec:eval-c1}

\argref{C1} asserts that mDNS/DNS-SD can provision AI services without end-user configuration. I decompose \argref{C1} into three requirements that together constitute an existence proof of feasibility, tracing the lifecycle of a zero-configuration AI connection. A service must first be \emph{found} without user input (R1). The discovery must then yield \emph{enough information} to connect (R2). Different implementations must agree on what they found, without sharing code (R3). If a service can be discovered, connected to, and consumed by independent clients, then zero-configuration provisioning is feasible.

\subsection{Results}

\paragraph{R1:}
Six implementations discover Saturn services through standard mDNS and DNS-SD queries on the \texttt{\_saturn.\_tcp.local.} service type. None requires a user-supplied URL, API key, or configuration file (\argref{A3}). The Python package, Rust router (Section~\ref{sec:router-implementation}), VLC extension, Open WebUI plugin, MCP server, and AI SDK provider each use their own mDNS library (\argref{A7}). The protocol, not my code, is what they share. This result connects directly to the access equity motivation from Chapter~\ref{ch:introduction}: if discovery requires no credentials and no configuration, then anyone on the network can use AI services regardless of whether they have a cloud provider account or a credit card.

\paragraph{R2:}
The TXT record schema (Section~\ref{sec:txt-schema}) carries everything a client needs to connect: endpoint URL, backend protocol, service priority, and hosting context (\argref{A8}). I tested this sufficiency against three backends with fundamentally different authentication models: Ollama requires none, DeepInfra uses a static JWT, and OpenRouter rotates ephemeral keys (\argref{A9}). All three work through the same schema with no out-of-band configuration necessary. This result matters beyond the protocol: schema sufficiency means any developer can integrate AI into an application without learning a provider's authentication model, billing API, or SDK. The schema carries that knowledge so the developer does not have to.

\paragraph{R3}
Seven consumers span three languages and four mDNS libraries (\argref{A10}): Python's \texttt{zeroconf}, Rust's \texttt{mdns-sd}, TypeScript's \texttt{multicast-dns}, and the macOS \texttt{dns-sd} CLI. No shared Saturn-specific code bridges them (\argref{A2}). Each parses standard DNS-SD TXT records independently. This is the result I find most telling: interoperability emerges from the protocol specification, not from a reference implementation I could have designed to make the numbers look good.

%----------------------------------------------------------------------
\section{C2: Reduced Configuration Effort}
\label{sec:eval-c2}

\argref{C2} asserts that Saturn reduces total configuration effort compared to per-user manual setup. I evaluate \argref{C2} using a cognitive walkthrough~\cite{wharton1994walkthrough}, counting discrete configuration actions each persona must perform to reach a single completion criterion: \emph{an end user on the network uses an AI-powered feature in an application.}

I define three personas that reflect Saturn's role model:

\begin{itemize}
  \item \textbf{Sysadmin} --- provisions the AI service and makes it available on the network.
  \item \textbf{App Developer} --- integrates AI capabilities into an application.
  \item \textbf{End User} --- consumes AI features through the developer's application.
\end{itemize}

Two conditions are compared. In the \emph{traditional} condition, the sysadmin obtains an API key from a cloud provider and distributes it to developers, who configure their applications with the key and expose AI features to end users, who create accounts and provide payment credentials. In the \emph{Saturn} condition, the sysadmin runs a Saturn beacon that advertises the service via mDNS; developers call \texttt{discover()} instead of configuring keys and URLs; end users do nothing.

The walkthrough scripts are implemented as annotated Python files in the project repository (\texttt{cog\_walkthrough/}), where each configuration action corresponds to a documented step. Step-counting optimizes for reproducibility and structural comparison: any reviewer following the same scripts arrives at the same counts. It sacrifices time measurement and cognitive difficulty weighting --- a limitation I address in Section~\ref{sec:eval-threats}. I chose a cognitive walkthrough over a user study because the relevant comparison is structural and because Saturn has no deployed user base from which to recruit participants.

\subsection{Results}

Table~\ref{tab:eval-cog} reports step counts per persona under each condition (\argref{A11}).

\begin{table}[ht]
\centering
\caption{Cognitive Walkthrough Step Counts by Persona}
\label{tab:eval-cog}
\begin{tabular}{l r r r}
\hline
\textbf{Persona} & \textbf{Traditional} & \textbf{Saturn} & \textbf{Change} \\
\hline
Sysadmin & 12 & 14 & +2 (+17\%) \\
App Developer & 19 & 4 & $-$15 ($-$79\%) \\
End User & 7 & 0 & $-$7 ($-$100\%) \\
\hline
\textbf{Total (1+1+1)} & \textbf{38} & \textbf{18} & $-$\textbf{20 ($-$53\%)} \\
\hline
\end{tabular}
\end{table}

The sysadmin does slightly more work: two additional steps to run a beacon process alongside the API key setup (\argref{CA3}). That is the cost. The payoff is stark. The app developer drops from nineteen steps to four, a 79\% reduction, because Saturn eliminates billing integration, key management, and endpoint configuration entirely (\argref{A12}). The end user drops to zero. No account creation, no payment credentials, no setup of any kind (\argref{A13}).

With one persona of each role, total steps drop from 38 to 18, a 53\% reduction. The reduction grows with scale because Saturn's formula has no $M$ term --- end users perform zero steps regardless of count (\argref{A14}):

\begin{align*}
\text{Traditional} &= 12 + 19N + 7M \\
\text{Saturn} &= 14 + 4N + 0M
\end{align*}

At $N = 10$ developers and $M = 100$ end users, this yields a 94\% reduction (902 vs.\ 54 steps).

The reduction is structural, not incremental: Saturn eliminates the entire billing and payment integration stack because the sysadmin's beacon already carries the credentials. Complexity centralizes in one administrator rather than scattering across $N$ developers and $M$ users. That is the core of \argref{C2}, and the step counts bear it out: 53\% reduction at minimum scale, 94\% at moderate scale. The reduction is not free --- the sysadmin absorbs two additional steps. But step counts treat all steps as equal, which they are not. Creating an API provider account is harder than running \texttt{pip install saturn}. I return to this limitation in Section~\ref{sec:eval-threats}.

%----------------------------------------------------------------------
\section{Threats to Validity}
\label{sec:eval-threats}

Three methodological threats constrain these conclusions. First, the three-persona model simplifies organizational reality. Roles overlap: a developer who also administers infrastructure would alter the per-persona step counts, though the structural elimination of billing and credential management persists regardless of role boundaries. Second, the cognitive walkthrough counts discrete actions but does not measure wall-clock time, error rates, or subjective difficulty. The identified steps --- package installation, service discovery, API calls --- are structurally simpler than traditional provisioning, but the methodology cannot confirm that they require less cognitive effort. Third, I, a single author, developed all seven service consumers and performed the walkthrough and the heuristic inspection reported in Section~\ref{sec:usability-critique}. Familiarity with Saturn’s architecture may consolidate multi-step processes, undercounting the friction a naive user would encounter. The executable Python walkthrough scripts provide a baseline for independent review, but single-author bias remains the most significant validity threat. Independent replication by external developers is required to robustly validate the feasibility and usability claims.

